%% 
\subsection{\deu{Statistisches Schließen}\eng{Statistical reasoning}}

\chapterIntro{%%
Einseitige Hypothesentests an der Binomialverteilung
}{%%
Voraussetzung ist ein Taschenrechner, der die Binomialverteilung kumuliert für mehrere Werte Gleichzeitig berechnen kann.

\begin{itemize}
    \item Die Beispiele wurden mit TI 30 X Plus MathPrint berechnet.
    \item Allen Beispielen liegt die Binomialverteilung zu Grunde.
    \item Es sollen Säulendiagramme zu den Verteilungen gezeichnet werden können.
    \item Bei gegebenem $\alpha$-Signifikanzniveau soll der Annahmebereich ($0..k$) und der Ablehnugsbereich ($k+1..n$) ermittelt werden können.
    \item Die Wahrscheinlichkeit für Fehler 1. und 2. Art ($\alpha$- und $\beta$-Fehler) werden verlangt
\end{itemize}
    
}%% end chapterIntro



\kNiveauAufgabe{
Die Statistiker eines Landes behaupten, die Anzahl weißer Personenwagen (PW) betrage \textbf{minimal} 20\,\%.
Ein Student hat darauf hin 260 PW ausgezählt und glaubt, es sind weniger als die 20\,\%. Die Nullhypothese ist $H_0: p=0.20$. 

a) Wie lautet die Alternativhypothese? 

b) Berechnen Sie den Ablehnungsbereich auf dem Signifikanzniveau $\alpha < 1\,\%$.
}{%% lösungen
a) $H_1$: Es sind weniger als 20\% der Autos weiß

b) Bei $k$ = 36 liefert binomial\textbf{c}df 0.637\,\% (danach bereits > 1\,\%). Der Annahmebereich "minimal 20" ist also [37;260]. Der Ablehnungsbereich liegt bei «nur 0 bis 36» Autos sind weiß.
}


%%%%%%%%%%%%%%%%%%%%%%%%%%%%%%%%%%%%%%%%%%%%%%%%%%%%%%%%%%%%%%%%%%%%%%%%%%%%%

\kNiveauAufgabe{
Christian und Philipp besitzen je einen klassischen sechsseitigen Spielwürfel. Ein Würfel wird 60-mal geworfen. Sei $X =$ Die Anzahl der geworfenen Sechsen.

a) Wir nehmen vorerst an, die beiden Würfel seien fair. Bestimmen Sie den Erwartungswert für die Anzahl der geworfenen Sechsen bei 60 Würfen. 

b) Bestimmen Sie die Wahrscheinlichkeit, dass in 60 Würfen mit einem fairen Würfel genau 14 Sechsen fallen.

\vspace{3mm}

Nun besteht im Laufe eines Spieleabends der Verdacht, dass der Würfel von Philipp manipuliert ist und häufiger eine 6 zeigt als ein fairer Würfel.

c) Formulieren Sie eine geeignete Nullhypothese H${}_0$ und eine Alternativhypothese H${}_1$ für diese Vermutung.

d) Bestimmen Sie auf dem Signifikanzniveau $\alpha = 1$\,\% eine Entscheidungsregel für eine Stichprobe von 60 Würfen, mit der beurteilt werden kann, ob der Würfel als fair angesehen wird.

e) Wie gross ist die Wahrscheinlichkeit für Fehler erster Art für die Entscheidungsregel aus Aufgabe d)? Was bedeutet ein Fehler erster Art am Beispiel mit dem Würfel?

f) Im Laufe des Abends gibt Philipp zu, dass sein Würfel tatsächlich manipuliert ist und mit $p = 0.3 =\frac{3}{10} $ die Augenzahl 6 erscheint. Berechnen Sie die Wahrscheinlichkeit für Fehler zweiter Art mit Ihrer Entscheidungsregel von Aufgabe d).
}{%% Lösungen
a) $$\mathbb{E}(\epsdice{6})=\frac{1}{6}\cdot{}60 = 10$$

b) 5.04\,\% = Binomialpdf(60, 1/6, 14)

c) $H_0$: Der Würfel ist fair und es fällt mit $\frac16$ Wahrscheinlichkeit eine \epsdice{6}.

$H_1$: Alternativhypothese: Der Würfel ist unfair und es fallen mit mehr als $\frac16$ Wahrscheinlichkeit eine \epsdice{6}.

d) Binomial\textbf{c}df(60, 1/6, 16) ist 98.35\,\%; Binomial\textbf{c}df(60,1/6, 17) ist hingegen 99.25\,\%. Somit ist von 0 bis und mit 17 mal eine \epsdice{6} zu werfen noch «normal» im $\alpha$=1\,\%-Signifikanzniveau. Die Entscheidungsregel lautet: «Der Annahmebereich der Nullhypothese liegt von 0 bis 17 mal eine \epsdice{6} zu werfen bei 60 mal werfen auf einem 1\,\% Signifikanzniveau.

e) Der Fehler 1. Art sagt, dass der faire Würfel als unfair taxiert wird. Bei einem Signifikanzniveau von 1\,\% müsste das auch 1\,\% sein. Da wir aber den Annahmebereich inklusive die 17 mal Werfen genommen haben, ist der Ablehnungsbereich = [18,60], was mit einer Wahrscheinlichkeit von 0.744\,\% eintreten wird: $P(X\ge 18) = 1-P(X\le 17) \approx 1- 0.9926 = 0.00744$

f) $P(X\le 17; 0.3; 60) \approx 45.14\,\%$ Mit ca 45\,\% Wahrscheinlichkeit wird ein gezinkter Würfel immer noch als fair angesehen, wenn wir [0;17] als Annahmebereich der Nullhypothese beibehalten.
}%

\kNiveauAufgabe{
Bei einer Lieferung von 250 Gläsern mit Brotaufstrich werden einige gefunden, die
zu wenig Gewicht haben. Der Vertrieb behauptet, dass mindestens 95\,\% das Norm-
gewicht haben. Zur Prüfung wird ein
linksseitiger Test mit 60 Gläsern und einem
Signifikanzniveau von 10\,\% durchgeführt;
49 Gläser haben das angegebene Gewicht, 11 nicht.

Gib den Ablehnungsbereich an und überprüfe das Ergebnis.
}{%% Lösungen
not yet
}%

\kNiveauAufgabe{
Bestimme mit einem linksseitigen Hypothesentest den Ablehnungsbereich
einer Stichprobe mit $n = 300$ und der Nullhypothese H${}_0$: $p = 0.4$ bei
einem Signifikanzniveau von

a) $\alpha = 0.1$

b) $\alpha = 0.05$

c) $\alpha = 0.01$

}{% Lösungen
}


\kNiveauAufgabe{
Laut einer Studie sind ca. 14\,\% der Bevölkerung in Deutschland Linkshänder. Berechne, wie groß
eine Gruppe zufällig ausgewählter Testpersonen mindestens sein muss, damit mit einer Wahrscheinlichkeit von mindestens 75\,\% mindestens

a) ein Linkshänder darunter ist, 

b) sechs Linkshänder darunter sind
}{% lösung
not yet
}


\kNiveauAufgabe{
Bei der Produktion von einfachen USB-Sticks als Werbegeschenke
sind erfahrungsgemäß einige defekt. Wie groß darf der durchschnittliche Anteil defekter
USB-Sticks höchstens sein, damit in einer Packung von 
200 USB-Sticks mit einer Wahrscheinlichkeit von 95\,\% höchstens 7 defekt sind?

Wählen Sie aus:

\begin{itemize}
    \item 1\,\%
    \item 1.5\,\%
    \item 2\,\%
    \item 2.5\,\%
\end{itemize}
}{%% Loesung
Ausprobieren: $p$ = $p$-Defekt $\approx$ 2.0057\,\%. Dabei habe ich maximal 7 defekte mit Wahrscheinlichkeit von 94.9997\,\%. Die beste Lösung aus der Auswahl sind 2\,\%. 
}

\kNiveauAufgabe{%% aufgabe
Marktforscher X.  behauptet, dass in 46 \% der Haushalte ein Mikrowellengerät
vorhanden ist; andere Marktforscher bezweifeln diese Angabe. Bei einer
Stichprobe des Statistischen Bundesamtes in 380 Haushalten wird
u.\,a. auch nach der Ausstattung des Haushalts gefragt.

a) Bei welchen Stichprobenergebnissen würde man die Angaben des Marktforschers 
X. für falsch ansehen?

Geben Sie  eine Entscheidungsregel auf dem $\alpha< 5\,\%$-Niveau  an!
 
b) Wie groß ist die Wahrscheinlichkeit für einen Fehler I. Art?
}{% Lösungen
not yet
}

%% %%%%%%%%%%%%%%%%%%%%%%% Aufgabe %%%%%%%%%%%%%%%%%%%%%%%%%%%%%%%%%%%%%%%%%%%%%%%%
\newpage
